\section{The Diagnosis: Coverage Prerequisite, Temporal Binding}
\label{sec:diag}

\subsection{A monotone correlation across three benchmarks}
The planner's gain over random policy selection ($-0.1$ / $+0.9$ / $+4.3$ on RVS / OVO / StreamingBench) tracks mean window coverage (0.21 / 0.59 / 1.00) almost perfectly — monotone in the thesis direction on both the full-question and all-wrong populations (correlational, three points; reported as a correlation, not a law).

\subsection{The ceiling decomposition}
Partitioning every question by whether any grain policy answers it correctly: RVS splits 110 all-right / 54 mixed / 152 all-wrong (48\%); StreamingBench 113 / 25 / 42 (23\%). Decomposing the all-wrong set: \textbf{zero retrieval failures} (no case where the evidence was present and unretrieved), 14 store-coverage failures, and a dominant model-capability remainder — with window-coverage evidence (mean 0.21 vs 0.071 chance; 83/138 temporal-before/after) localizing the capability failure to \emph{long-window temporal grounding}: the model cannot integrate evidence across a wide temporal span even when the store holds it. StreamingBench's all-wrong set is capability-failure by construction (its windows are short and end at the stream frontier, where coverage is 1.0).

\subsection{The causal test: refuted on evidence}
We attempted to force the diagnosis: four allocation policies (top-k, uniform-spread, hybrid, stratified) over one store and one 64-slot budget, spread computed strictly over streamed history (never the gold window — asserted in code), on the 83 temporal questions, with a pre-registered decision mapping. The deciding two-arm test was underpowered ($p{=}0.180$ vs a pre-declared MDE of 16--21 points); we therefore ran the powered version: a pooled regression over all 332 (question, arm) observations with question fixed effects and cluster-robust errors. \textbf{The result is a refutation, not a null-of-power:} the coverage coefficient is flat and insignificant in every specification (question-FE: $-0.33$ per unit coverage, CI $[-1.99, +1.33]$), the coverage--accuracy curve is non-monotone across six bins, and the one arm with a real accuracy gain (hybrid, $+6$ pt) had \emph{lower} coverage than the baseline. Within the achieved coverage band ($\approx$0.18--0.49 — the arms spread retrieval \emph{content}, Jaccard 0.18--0.29, but barely coverage, cross-arm correlation 0.94--0.99), \textbf{coverage is not the lever}. Combined with the oracle-null result, the closure is complete: on this system, these questions are answered wrong regardless of which evidence the allocator selects — the binding constraint is model capability at long temporal windows, and no allocation mechanism at this store budget reaches the coverage where that might change.

\subsection{What the field is optimizing}
Grain selection, bit-width allocation, and retrieval ordering — the literature's active axes — are second-order terms twice over: the oracle headroom they chase is an artifact (Section~\ref{sec:null}), and the coverage lever that might have mattered is refuted (this section). The binding constraint is the model's long-window temporal grounding, which current streaming benchmarks' standard splits barely probe and current mechanisms do not address. That is where an accuracy claim could live, and it is out of scope for a measurement paper.
